\documentclass[12pt]{extarticle}

\usepackage{tabularx}
\usepackage{amsmath}
\usepackage{extsizes}
\usepackage{graphicx}
\usepackage{amssymb}
\usepackage[english,russian]{babel}
\usepackage[utf8]{inputenc}
\usepackage[a4paper]{geometry}
\geometry{top=20mm,bottom=30mm,left=20mm,right=15mm}
\usepackage[T1,T2A]{fontenc}
\usepackage{setspace}
\usepackage{indentfirst}
\usepackage{caption}
\renewcommand{\baselinestretch}{1.5}
\usepackage{spverbatim}
\usepackage{listings}
\usepackage{xcolor}

\definecolor{codebg}{gray}{0.96}
\definecolor{comment}{gray}{0.45}
\definecolor{keyword}{RGB}{0, 75, 150}
\definecolor{string}{RGB}{150, 80, 30}
\definecolor{number}{gray}{0.5}
\definecolor{identifier}{rgb}{0,0,0}

\lstdefinestyle{coolprint}{
  basicstyle   = \ttfamily\footnotesize,
  keywordstyle = \bfseries\color{keyword},
  commentstyle = \itshape\color{comment},
  stringstyle  = \color{string},
  identifierstyle = \color{identifier},
  backgroundcolor = \color{codebg},
  numbers      = left,
  numberstyle  = \footnotesize\color{number},
  numbersep    = 5pt,
  frame        = single,
  framerule    = 0.8pt,
  framesep     = 3pt,
  breaklines   = true,
  breakatwhitespace = true,
  tabsize      = 4,
  showspaces   = false,
  showstringspaces = false,
  showtabs     = false,
  xleftmargin  = 10pt,
  xrightmargin = 5pt,
  aboveskip    = \medskipamount,
  belowskip    = \medskipamount,
  captionpos   = b,
}
\lstset{style=coolprint}

\begin{document}

\begin{center}
  \textbf{МИНОБРНАУКИ РОССИИ} \\
  \textbf{Санкт-Петербургский государственный} \\
  \textbf{электротехнический университет} \\
  \textbf{«ЛЭТИ» им. В.И. Ульянова (Ленина)} \\
  \textbf{Кафедра САПР} \\
\end{center}

\vspace{0.25\textheight}

\begin{center}
\textbf{ОТЧЕТ} \\
\textbf{по лабораторной работе №3 (9)} \\
\textbf{по дисциплине «Автобусы»} \\
\textbf{Тема: КРИТИЧЕСКИЕ ИСПЫТАНИЕ И КРАШТЕСТ} \\
\end{center}

\vspace*{\fill}

\begin{center}
\begin{tabularx}{\textwidth}{ l X r }
 Студенты гр. 5352 & \rule{8cm}{0.2mm} & Cоболевский М.А. \\
 & \rule{8cm}{0.2mm} & Решетников С.А. \\
 Преподаватель & \rule{8cm}{0.2mm} & Кочетков А.В. \\
\end{tabularx}
\end{center}

\vspace{3em}

\begin{center}
Санкт-Петербург\\
\the\year
\end{center}

\thispagestyle{empty}

\pagebreak

\section{Задание на лабораторную работу}

Программа использует два целых значения $x$ и $y$ из секции \texttt{.data}, вычисляет
следующее выражение $(x >> 2) + ((y - 1) << 3)$ и выводит результат.
($>>$ - логический сдвиг)

\section{Текст программы}

\renewcommand{\baselinestretch}{1}
\input{assets/l1}
\renewcommand{\baselinestretch}{1.5}

\section{Примеры запуска}

\section{Выводы}

\end{document}
